\begin{figure}[htbp]
\centering
\begin{tikzpicture}[x=.95cm,y=.95cm,font=\footnotesize,>={Stealth[length=1.8mm]}]
\draw[red,semithick] (2.8,2.9)--(2.8,5.5);
\draw[blue!65!cyan,semithick] (2.8,2.9)--(5.5,2.9);
\draw[green!65!black,semithick] (2.8,5.5)--(5.5,2.9);
\node[above] at (2.8,5.5) {$X$};
\node[right] at (5.5,2.9) {$Y$};
\node[anchor=east] at (2.6,4.15) {$\lVert X\rVert^2=\Var(X)$};
\node[anchor=west] at (4.2,4.8) {$\lVert X+Y\rVert^2=\Var(X+Y)$};
\node[below] at (4.15,2.5) {$\lVert Y\rVert^2=\Var(Y)$};
\node at (3.7,1.25) {$\Var(X+Y)=\Var(X)+\Var(Y)$};
\node at (.5,.2) {（a）};
\draw[red!75!black,dash dot,semithick] (8,1.15)--(14,1.15)--(14.9,3.6)--(8.9,3.6)--cycle;
\node at (8.9,1.65) {$E$};
\draw[red,semithick] (10.4,6.1) .. controls (11.25,5.9) and (11.4,2.25) .. (10.1,3.25);
\draw[green!65!black,semithick] (10.4,6.1) .. controls (12.4,7.15) and (13.85,5.55) .. (13.65,3.42);
\draw[blue!65!cyan,dashed,semithick] (10.8,3.15) .. controls (12.0,3.10) and (13.45,3.05) .. (13.65,3.42);
\filldraw[fill=black,draw=inkblue,thick] (10.4,6.1) circle (.09);
\filldraw[fill=black,draw=inkblue,thick] (10.1,3.25) circle (.09);
\filldraw[fill=black,draw=inkblue,thick] (13.65,3.42) circle (.09);
\node[left] at (10.3,6.1) {$P(x)$};
\node[left] at (10.0,3.25) {$U(x)$};
\node[right] at (13.6,3.05) {$V_g(x)$};
\node at (9.8,4.75) {$D(P\Vert U)$};
\node at (12.2,4.45) {$D(P\Vert V_g)$};
\node at (12.0,2.78) {$D(U\Vert V_g)$};
\draw[->,dashed,magenta!65!gray,semithick] (10.0,3.02)--(10.55,2.0);
\draw[->,dashed,magenta!65!gray,semithick] (12.5,1.98)--(13.6,3.26);
\node[below] at (10.55,2.0) {$U_f(y)$};
\node[below] at (12.5,1.98) {$V(y)$};
\node at (11.5,.4) {$D(P\Vert V_g)=D(P\Vert U)+D(U\Vert V_g)$};
\node at (7.8,.2) {（b）};
\end{tikzpicture}
\origfig{1}{勾股定理。（a）勾股定理与协方差；（b）勾股定理与独立性。}
\label{fig:ch07a-01}
\end{figure}
